% p2-01-motivation

\section{P2 §1. Motivation and Context}

1. Motivation and Context

      Figure (fig-p2-open). Opening triptych --- Ladder / Wheel / Mirror. Pellis additive expansion vs
                   Trinity multiplicative monomials, joined at the isomorphism log phi.

  1.1 What Paper 1 Established
  The first paper in this programme performed a symbolic compressibility analysis on a
  curated set of dimensionless combinations of Standard Model parameters and related
  quantum critical observables. The central finding was that a subset of these

  combinations --- particularly those involving gauge coupling ratios near threshold
  energies, Yukawa coupling hierarchy indicators, and measured mass ratios from
  near-critical condensed matter systems --- can be expressed as short algebraic words
  over the alphabet {$\varphi$, pi, e, 2, 3, 5, +, -, $\times$, /} more economically than over comparison
  alphabets that exclude $\varphi$. The compression advantage, measured by the ratio of
  minimal description lengths, was statistically significant in the low-to-moderate
  complexity tier but not universal.

  Paper 1 made no mechanistic claim. Its conclusion was: the observed $\varphi$-compressibility
  is a fact about symbolic structure, and a mechanism explanation is warranted only if (a)
  the effect is robust to sampling and encoding choices, and (b) there exists at least one
  identified pathway by which physical dynamics could produce it. The present paper
  addresses condition (b).

  Relationship to the PhD parameter catalogue. The Flos Aureus monograph (Zenodo
  10.5281/zenodo.19227877) maintains a 42-parameter catalogue (Chapter flos34) that
  includes coupling constant combinations with strong fit-density claims and Monte Carlo
  support. That catalogue is a candidate data source for the Paper 1 statistical filtering
  pass; however, it has not yet been subjected to the tolerance-variation and
  encoding-independence tests that Paper 1 requires. In the present paper, flos34 is cited
  as a catalogue requiring Paper 1 statistical filtering --- not as independent empirical
  confirmation of $\varphi$-structured couplings. The filtering is Work Package 5.

  1.2 The Mechanistic Question
  A $\varphi$-pattern in physical constants could arise through several non-exclusive channels:

  1. Algebraic coincidence. The golden ratio satisfies $\varphi$2 = $\varphi$ + 1 and therefore appears
  frequently in algebraic expressions over Z[$\varphi$]. Any sufficiently algebraic number system
  will encounter $\varphi$ without physical cause.

  2. Symmetry-group structure. Certain Lie algebras and Coxeter groups have
  Perron-Frobenius eigenvectors whose components are $\varphi$-valued. If a physical system's
  mass spectrum is governed by such an algebra, $\varphi$ appears as an exact kinematic
  consequence --- as in the Zamolodchikov E8 case (Project Euclid, Adv. Stud. Pure Math.
  19 (1989)).

  3. Renormalization group fixed-point geometry. If RG flows possess fixed points or limit
  cycles with self-similar structure under scale transformations by ratio lambda = $\varphi$, then
  $\varphi$-structured ratios would be dynamically preferred near those fixed points (Wilson and
  Kogut, Phys. Rep. 12 (1974) 75).

  4. Discrete scale invariance. Systems whose symmetry group is a discrete subgroup of
  the conformal group rather than the full conformal group exhibit log-periodic
  corrections to scaling with a preferred ratio lambda. If lambda = $\varphi$, $\varphi$ enters the
  spectrum of observables (Sornette, Phys. Rep. 297 (1998) 239).

  This paper examines whether channels 2, 3, and 4 have concrete field-theoretic
  realizations, and whether those realizations connect to the patterns observed in Paper
  1. Channel 1 is kept as the null hypothesis throughout.

  1.3 The AlphaPhi Notation Boundary

  The Flos Aureus monograph Chapter flos38 defines candidate $\alpha$-$\varphi$ expressions. An
  internal notation inconsistency must not be inherited by this paper. The candidates in
  that chapter are structurally distinct and are here given permanent, unambiguous
  names for use throughout Papers 2 and 3:

  alpha$\varphi$,log and alpha$\varphi$,alg are numerically and structurally distinct; they must never be
  conflated. Neither is derived from first principles in the present paper. The
  disambiguation of these three quantities is a prerequisite for Paper 3. All coupling
  constant expressions in the present paper define their notation explicitly upon first use.